\documentclass[10pt, a4paper]{article}
\usepackage[paper=a4paper, left=1.5cm, right=1.5cm, bottom=1.5cm, top=3.5cm]{geometry}
\usepackage[utf8]{inputenc}
% \usepackage[T1]{fontenc}
\usepackage[spanish]{babel}
\usepackage{indentfirst}
\usepackage{fancyhdr}
\usepackage{latexsym}
\usepackage{lastpage}
\usepackage{framed}
\usepackage{todonotes} % para dejar notitas de to-do!
\usepackage{aed2-symb,aed2-itef,aed2-tad,aed2-diseno}
\usepackage[colorlinks=true, linkcolor=black]{hyperref}
\usepackage{calc}
\usepackage{graphicx}
\usepackage{changepage}
%

% ========== Para escribir pseudo ==========
\usepackage{algorithm}
\usepackage[noend]{algpseudocode}  % "noend" es para no mostrar los endfor, endif
\algrenewcommand\alglinenumber[1]{\tiny #1:}  % Para que los numeros de linea del pseudo sean pequeños
\renewcommand{\thealgorithm}{}  % Que no aparezca el numero luego de "Algorithm"
\floatname{algorithm}{ }    % Entre {  } que quiero que aparezca en vez de "Algorithm"

% traducciones
% \algrenewcommand\algorithmicwhile{\textbf{mientras}}
% \algrenewcommand\algorithmicdo{\textbf{hacer}}
% \algrenewcommand\algorithmicreturn{\textbf{devolver}}
% \algrenewcommand\algorithmicif{\textbf{si}}
% \algrenewcommand\algorithmicthen{\textbf{entonces}}
% \algrenewcommand\algorithmicfor{\textbf{para}}
%=========================================================


\newcommand{\f}[1]{\text{#1}}
\renewcommand{\paratodo}[2]{$\forall~#2$: #1}
\newcommand{\numeroEjercicio}[1]{\textbf{\large{Ejercicio #1:}}\\}
\newcommand{\tituloSubEjercicio}[1]{$\newline$\tadNombre{#1:}}

\sloppy

\hypersetup{%
 % Para que el PDF se abra a página completa.
 pdfstartview= {FitH \hypercalcbp{\paperheight-\topmargin-1in-\headheight}},
 pdfauthor={Cátedra de Algoritmos y Estructuras de Datos II - DC - UBA},
 pdfkeywords={Template TADs básicos},
 pdfsubject={Template TADs básicos}
}

\parskip=5pt % 10pt es el tamaño de fuente

% Pongo en 0 la distancia extra entre ítemes.
\let\olditemize\itemize
\def\itemize{\olditemize\itemsep=0pt}

% Acomodo fancyhdr.
\pagestyle{fancy}
\thispagestyle{fancy}
\addtolength{\headheight}{1pt}
\lhead{Algoritmos y Estructuras de Datos II}
\rhead{$2^{\mathrm{do}}$ cuatrimestre de 2022}
\cfoot{\thepage /\pageref{LastPage}}
\renewcommand{\footrulewidth}{0.4pt}

\author{Algoritmos y Estructuras de Datos II, DC, UBA.}
\date{}
\title{Trabajo Práctico de Algoritmos y Estructuras de Datos II}

\NeedsTeXFormat{LaTeX2e}
\ProvidesPackage{caratula}[2003/4/13 v0.1 Para componer caratulas de TPs del DC]


% ----- Imprimir un mensajito al procesar un .tex que use este package -----

\typeout{Cargando package `caratula' v0.2 (21/4/2003)}


% ----- Algunas variables --------------------------------------------------

\let\Materia\relax
\let\Submateria\relax
\let\Titulo\relax
\let\Subtitulo\relax
\let\Grupo\relax


% ----- Comandos para que el usuario defina las variables ------------------

\def\materia#1{\def\Materia{#1}}
\def\submateria#1{\def\Submateria{#1}}
\def\titulo#1{\def\Titulo{#1}}
\def\subtitulo#1{\def\Subtitulo{#1}}
\def\grupo#1{\def\Grupo{#1}}


% ----- Token list para los integrantes ------------------------------------

\newtoks\intlist\intlist={}


% ----- Comando para que el usuario agregue integrantes

\def\integrante#1#2#3{\intlist=\expandafter{\the\intlist
    \rule{0pt}{1.2em}#1&#2&\tt #3\\[0.2em]}}


% ----- Macro para generar la tabla de integrantes -------------------------

\integrante{Guberman, Diego Andrés}{469/17}{diego98g@hotmail.com}
\integrante{Ramis Folberg, Ezequiel Leonel}{881/21}{ezequielramis.hello@gmail.com}
\integrante{Sabetay, Kevin Damian}{476/16}{kevin.sabetay96@gmail.com}

\def\tablaints{
    \begin{tabular}{|l@{\hspace{4ex}}c@{\hspace{4ex}}l|}
        \hline
        \rule{0pt}{1.2em}Integrante & LU & Correo electr\'onico\\[0.2em]
        \hline
        \the\intlist
        \hline
    \end{tabular}}

% ----- Macro para generar la parte reservada para la c�tedra -------------------------

\def\tablacatedra{%
    \\
    \textbf{Reservado para la c\'atedra}\par\bigskip
    \begin{tabular}{|c|c|c|}
        \hline
        \rule{0pt}{1.2em}Instancia & Docente & Nota\\[0.2em]
        \hline
        \rule{0pt}{1.2em}Primera entrega & \phantom{mmmmmmmmmmmmmmmmmm} & \phantom{mmmmmm} \\
        \hline
        \rule{0pt}{1.2em}Segunda entrega & & \\
        \hline
    \end{tabular}}

% ----- Codigo para manejo de errores --------------------------------------

\def\se{\let\ifsetuperror\iftrue}
\def\ifsetuperror{%
    \let\ifsetuperror\iffalse
    \ifx\Materia\relax\se\errhelp={Te olvidaste de proveer una \materia{}.}\fi
    \ifx\Titulo\relax\se\errhelp={Te olvidaste de proveer un \titulo{}.}\fi
    \edef\mlist{\the\intlist}\ifx\mlist\empty\se%
    \errhelp={Tenes que proveer al menos un \integrante{nombre}{lu}{email}.}\fi
    \expandafter\ifsetuperror}


% ----- Reemplazamos el comando \maketitle de LaTeX con el nuestro ---------

\def\maketitle{%
    \thispagestyle{empty}
    \begin{center}
    \vspace*{\stretch{2}}
    \materia{Algoritmos y Estructuras de Datos II}

    {\LARGE\textbf{\Materia}}\\[1em]
    \submateria{Trabajo Pr\'actico 1}
    \ifx\Submateria\relax\else{\Large \Submateria}\\[0.5em]\fi

%\def\titulo#1{\def\Titulo{#1}}
%\def\subtitulo#1{\def\Subtitulo{#1}}
%\def\grupo#1{\def\Grupo{#1}}
    \par\vspace{\stretch{1}}
    \begin{figure}[h!]
      \centering
      \includegraphics[width=0.5\linewidth]{caratula.png}
    \end{figure}
    \par\vspace{\stretch{1}}
    {\large Departamento de Computaci\'on}\\[0.5em]
    {\large Facultad de Ciencias Exactas y Naturales}\\[0.5em]
    {\large Universidad de Buenos Aires}
    \par\vspace{\stretch{3}}
    {\Large \textbf{\Titulo}}\\[0.8em]
    {\Large \Subtitulo}
    \par\vspace{\stretch{3}}
    \grupo
    \ifx\Grupo\relax\else\textbf{\Grupo}\par\bigskip\fi
    \tablaints
    \vspace*{\stretch{3}}
    \medskip
    \tablacatedra
    \end{center}
    \vspace*{\stretch{3}}
    \newpage
    }

% Comandos para cositas de complejidad

\newcommand{\bigO}{\mathcal{O}}
\newcommand{\Nat}{\mathbb{N}}
\newcommand{\R}{\mathbb{R}}
\newcommand{\Rpos}{\mathbb{R}_{>0}}
\newcommand{\eqdef}{\overset{\mathrm{def}}{=}}
\newcommand{\eqprop}{\overset{\mathrm{prop}}{=}}
%\newcommand{\ssi}{\leftrightarrow}

\newcommand{\seExplica}[2]{
\textbf{se explica con}: \NombreFuncion{#1}\par\noindent
\textbf{géneros}: \TipoVariable{#2}
}

\newcommand{\usa}[1] {
\textbf{usa}: \NombreFuncion{#1}\par
}

\newenvironment{interfaz}[1]
  {#1
  \adjustwidth{3em}{}}
  {\endadjustwidth}

\newenvironment{operaciones}
  {\textbf{operaciones}:
  \adjustwidth{2em}{}}
  {\endadjustwidth}

\newcommand{\asignar}[2]{$#1 \gets #2$}

\newcommand{\falta}{\todo[inline, inlinewidth=1.7em, noinlinepar]{??}}

\newcommand{\defun}[3]{
  \caption{\textsc{#1}(#2) $\longrightarrow$ \texttt{#3}}
}

\newcommand{\NULL}{\textsc{NULL}}
\newcommand{\Lmax}{L_{\texttt{máx}}}

\begin{document}

\maketitle

\tableofcontents
\newpage

\section{Dudas}
\begin{itemize}
        \item Igualdad observacional sin tenerla definida en los TADs.
        \item Nuevo Juego: ¿El repo se pasa por parámetro o se genera ``aleatoriamente'' dentro de la función como un comportamiento automático? Si es el 1º caso se tendría que copiar para coincidir con la complejidad del ejercicio. Si es el segundo ver cómo se genera ``aleatoriamente''.
        \item ¿Qué tipo es \texttt{letra}? Si es \texttt{char} la complejidad de nuevoJuego no se cumple ($\Sigma$ sería muy grande). En caso de pasar por parámetro un conjunto de \texttt{chars} como alfabeto, habría que checkear que todas las fichas del repo inicial sean de ese alfabeto. Lo más probable es que sea un \texttt{enum}.
        \item ¿Se puede usar operaciones de interfaz dentro de \textbf{Pre} y \textbf{Post}? Por ejemplo para acceder a elemento de tablero, si es otro módulo.
        \item Si se quiere que \texttt{variante} sea un módulo, ¿qué tipo de parámetro de entrada se usan para la operación ``nuevaVariante''? ¿Las ya definidas en la representación? ¿Puede ser un módulo que no tenga interfaz y solamente el \textbf{Rep} y \textbf{Abs}?
        \item ¿Hay que poner el item ``usa'' o sólo completar los ``requiere''?
        \item Complejidad de operación \textsc{\#letraTieneJugador}.
\end{itemize}

\section{Preámbulo}
Antes de presentar los módulos, definimos las siguientes variables para las complejidades temporales:
\begin{itemize}
    \item $N$ \---- tamaño del tablero.
    \item $K$ \---- cantidad de jugadores.
    \item $|\Sigma|$ \---- cantidad de letras en el alfabeto.
    \item $F$ \---- cantidad de fichas por jugador.
    \item $L_{\texttt{máx}}$ \---- longitud de la palabra legítima más larga definida por la variante del juego de la que se trate.
\end{itemize}

%%% Local Variables:
%%% mode: latex
%%% TeX-master: "main"
%%% End:

\section{Módulo Juego}
\begin{interfaz}{\subsection{Interfaz}}
  \seExplica{Juego}{juego}\\
  \usa{\falta}
  \par\noindent
  \begin{operaciones}
    \InterfazFuncion{nuevoJuego}
    {
      \In{k}{nat},
      \In{v}{variante}
    }{juego}
    [$k>0$]
    {$\exists r: \texttt{cola(letra)} \mid res \igobs nuevoJuego(k,v,r)$}
    [$O(N^{2}+|\Sigma|K+FK)$]
    [Dada una cantidad de jugadores y una variante de juego, se inicia un nuevo juego con el tablero vacío y con un repositorio de fichas acorde.]
    [\falta]
    [\falta]\\

    \noindent\InterfazFuncion{jugadaVálida?}
    {
      \In{j}{juego},
      \In{o}{occurrencia}
    }{bool}
    [true]
    {$res\igobs jugadaV\acute{a}lida?(j,o)$}
    [$O(L_{\texttt{máx}}^{2})$]
    [Determina si una jugada es válida.]
    [\falta]
    [\falta]\\

    \noindent\InterfazFuncion{ubicar}
    {
      \Inout{j}{juego},
      \In{o}{occurrencia}
    }{}
    [$jugadaV\acute{a}lida(j,o) \land j\igobs J_{0} $]
    {$j\igobs ubicar(J_{0},o)$}
    [$O(m)$, donde $m$ es el número de fichas que se ubican.]
    [Ubica un conjunto de fichas en el juego.]
    [\falta]
    [\falta]\\

    \noindent\InterfazFuncion{variante}
    {
      \In{j}{juego}
    }{variante}
    [true]
    {$res\igobs variante(j)$}
    [$O(1)$]
    [Obtiene información sobre la variante del juego.]
    [\falta]
    [\falta]\\

    \noindent\InterfazFuncion{turno}
    {
      \In{j}{juego}
    }{nat}
    [true]
    {$res\igobs turno(j)$}
    [$O(1)$]
    [Obtiene al jugador del turno actual.]
    [\falta]
    [\falta]\\

    \noindent\InterfazFuncion{puntaje}
    {
      \In{j}{juego},
      \In{i}{nat}
    }{nat}
    [$i < \#jugadores(j)$]
    {$res\igobs puntaje(j,i)$}
    [$O(m\cdot L_{\texttt{máx}})$, donde $m$ es la cantidad de fichas que ubicó el
jugador desde la última vez que se invocó a esta operación.]
    [Obtiene el puntaje de un jugador.]
    [\falta]
    [\falta]\\

    \noindent\InterfazFuncion{celda}
    {
      \In{J}{juego},
      \In{i}{nat},
      \In{j}{nat}
    }{puntero(letra)}
    [$enTablero?(tablero(J),i,j)$]
    {$res\igobs\IFL{hayLetra?(tablero(J),i,j)}THEN{\&letra(tablero(J),i,j)}ELSE{\NULL}FI$}
    [$O(1)$]
    [Obtiene el contenido del tablero en una coordenada $(i,j)$.]
    [\falta]
    [\falta]\\

    \noindent\InterfazFuncion{\#letraTieneJugador}
    {
      \In{j}{juego},
      \In{x}{letra},
      \In{i}{nat}
    }{nat}
    [$i < \#jugadores(j)$]
    {$res\igobs \#(x,fichas(j,i))$}
    [$O(1)$]
    [Dada una cierta letra $x$ del alfabeto, conocer cuántas fichas tiene un jugador de dicha letra.]
    [\falta]
    [\falta]\\

  \end{operaciones}
\end{interfaz}

\subsection{Implementación}
\subsubsection*{Representación}
\begin{Estructura}{foo}
\begin{Tupla}
  \tupItem{foo}{bar}
  \tupItem{foo}{bar}
\end{Tupla}
\end{Estructura}

\subsubsection*{Invariante de Representación}
\Rep{foo}

\par\vspace*{3ex}%

\subsubsection*{Función de Abstracción}
\Abs[estr]{foo}[e]{p}{bar}

\subsubsection*{Algoritmos}
\begin{algorithm}[H]
  \defun{hacerGuia}{\In{A}{guia}, \In{parámetroInútil}{Nat}}{bool}
  \begin{algorithmic}[1]
    \State\asignar{i}{0} \Comment{esto es $\Theta$(1)}
    \State\asignar{n}{guia.\text{cantEjercicios()}} \Comment{$\bigO$(1)}
    \State\asignar{consultas}{\textsc{diccVacio}}
    \State\textsc{prepararMate()} \Comment{$\Omega$($n^n$)}
    \While{$i < n$}
    \State\textsc{pensarEjercicio(i)}
    \If{\textsc{tengoConsultas($i$)}}
      \State\textsc{escribirConsultasEjercicio($i$,$consultas$)}
    \Else
      \State\textsc{comerBizochito()}
    \EndIf
      \State\textsc{comerBizochito()}
    \EndWhile
    \For{miVariable}
         \State hacer algo
    \EndFor
    \State\Return{\textsc{vacio?}($consultas$)}
  \end{algorithmic}
\end{algorithm}
\newpage

%%% Local Variables:
%%% mode: latex
%%% TeX-master: "main"
%%% End:

\section{Módulo Servidor}
\subsection{Interfaz}
\subsection{Implementación}
\newpage

%%% Local Variables:
%%% mode: latex
%%% TeX-master: "main"
%%% End:

\section{Módulos auxiliares}
\subsection{Módulo Letra}
Se asume una implementación acorde\footnote{Una buena opción es usar un \href{https://en.wikipedia.org/wiki/Enumerated_type}{Enumerado}.} al módulo de género \texttt{letra} con las siguientes operaciones en la interfaz (todas con órden de complejidad $O(1)$):
\begin{itemize}
    \item $\textsc{dom}\quad :\quad\quad\quad \rightarrow\texttt{nat}$ \---- Tamaño del dominio del tipo \texttt{letra}. Corresponde con la variable $A$ de su especificación.
    \item $\textsc{ord}\quad:\texttt{letra}\rightarrow\texttt{nat}$ \---- Dada una letra, devuelve su correspondiente índice.
    \item $\textsc{ord}^{-1}: \texttt{nat}\ n\rightarrow\texttt{letra}\ \{n < A\}$ \---- Dado un índice, devuelve su correspondiente letra.
\end{itemize}

\subsection{Módulo Variante (Trie)}
\begin{interfaz}{\subsubsection{Interfaz}}
  \seExplica{Variante}{variante}\\
  \usa{\falta}
  \par\noindent
  \begin{operaciones}
    \InterfazFuncion{nuevaVariante}
    {
      \begin{adjustwidth}{2em}{}
      \In{n}{nat},\\
      \In{f}{nat},\\
      \In{puntajes}{dicc(letra, nat)},\\
      \In{legítimas}{conj(secu(letra))}
      \end{adjustwidth}
    }{variante}
    [$n>0\land f>0$]
    {$res\igobs nuevaVariante(n,f,puntajes,leg\acute{\imath}timas)$}
    [$O(\#leg\acute{\imath}timas\cdot\Lmax)$]
    [Genera una variante de juego.]
    [\falta]\\

    \noindent\InterfazFuncion{tamañoTablero}
    {\In{v}{variante}}{nat}
    [true]
    {$res\igobs tama\tilde{n}oTablero(v)$}
    [$O(1)$]
    [Devuelve el tamaño del tablero.]
    [\falta]\\

    \noindent\InterfazFuncion{fichasPorJugador}
    {\In{v}{variante}}{nat}
    [true]
    {$res\igobs \#fichas(v)$}
    [$O(1)$]
    [Devuelve la cantidad de fichas que debe de tener cada jugador.]
    [\falta]\\

    \noindent\InterfazFuncion{puntajeLetra}
    {\In{v}{variante}, \In{l}{letra}}{nat}
    [true]
    {$res\igobs puntajeLetra(v,l)$}
    [$O(1)$]
    [Devuelve el puntaje de una letra.]
    [\falta]\\

    \noindent\InterfazFuncion{palabraLegítima?}
    {\In{v}{variante}, \In{l}{secu(letra)}}{bool}
    [true]
    {$res\igobs palabraLeg\acute{\imath}tima(v,l)$}
    [$O(\Lmax)$]
    [Determina si una palabra es legítima dentro de la variante de juego.]
    [\falta]\\

    \noindent\InterfazFuncion{longPalabraMásLarga}
    {\In{v}{variante}}{nat}
    [true]
    {
      \begin{adjustwidth}{2em}{}

      $(\exists p:\texttt{secu(letra)})(res\igobs long(p)\land palabraLeg\acute{\imath}tima?(v,p)\land\\(\forall p_{2}:\texttt{secu(letra)})( palabraLeg\acute{\imath}tima?(v,p_{2})\implies res\geq long(p_{2})))$
      \end{adjustwidth}
    }
    [$O(1)$]
    [Obtiene la longitud de la palabra legítima más larga de la variante.]
    [\falta]
  \end{operaciones}
\end{interfaz}
\subsubsection{Implementación}

\subsection{Módulo Tablero}
\begin{interfaz}{\subsubsection{Interfaz}}
  \seExplica{Tablero}{tab}\\
  \usa{\falta}
  \par\noindent
  \begin{operaciones}
    \InterfazFuncion{nuevoTablero}
    {
      \In{n}{nat}
    }{tab}
    [$n>0$]
    {$res\igobs nuevoTablero(n)$}
    [$O(N^{2})$]
    [Genera un tablero de tamaño $n$.]
    [\falta]\\

    \noindent\InterfazFuncion{tamaño}
    {
      \In{t}{tab}
    }{nat}
    [true]
    {$res\igobs tama\tilde{n}o(t)$}
    [$O(1)$]
    [Devuelve tamaño del tablero.]
    [\falta]\\

    \noindent\InterfazFuncion{ponerLetra}
    {
      \In{t}{tab},
      \In{i}{nat},
      \In{j}{nat},
      \In{l}{letra},
      \In{tm}{nat}
    }{tab}
    [$enTablero(t,i,j)\yluego\lnot hayLetra?(t,i,j)$]
    {$res\igobs ponerLetra(t,i,j,l)$}
    [$O(1)$]
    [Pone una letra en la coordenada $(i,j)$.]
    [\falta]\\

    \noindent\InterfazFuncion{sacarLetra}
    {
      \In{t}{tab},
      \In{i}{nat},
      \In{j}{nat}
    }{tab}
    [$enTablero(t,i,j)\yluego hayLetra?(t,i,j)$]
    {$res$ es igual al tablero antes de haber puesto una letra en la coordenada $(i,j)$}
    [$O(1)$]
    [Saca una letra en la coordenada $(i,j)$.]
    [\falta]\\

    \noindent\InterfazFuncion{enTablero?}
    {\In{t}{tab}, \In{i}{nat}, \In{j}{nat}}{bool}
    [true]
    {$res\igobs enTablero?(t,i,j)$}
    [$O(1)$]
    [Determina si una coordenada $(i,j)$ está en el rango del tablero.]
    [\falta]\\

    \noindent\InterfazFuncion{hayLetra?}
    {\In{t}{tab}, \In{i}{nat}, \In{j}{nat}}{bool}
    [$enTablero?(t,i,j)$]
    {$res\igobs hayLetra?(t,i,j)$}
    [$O(1)$]
    [Determina si una celda del tablero dada una coordenada $(i,j)$ está ocupada por una letra.]
    [\falta]\\

    \noindent\InterfazFuncion{letra}
    {\In{t}{tab}, \In{i}{nat}, \In{j}{nat}}{letra}
    [$enTablero?(t,i,j)\yluego hayLetra?(t,i,j)$]
    {$res\igobs letra(t,i,j)$}
    [$O(1)$]
    [Obtiene la letra de una ficha del tablero dada una coordenada $(i,j)$.]
    [\falta]\\

    \noindent\InterfazFuncion{tiempo}
    {\In{t}{tab}, \In{i}{nat}, \In{j}{nat}}{nat}
    [$enTablero?(t,i,j)\yluego hayLetra?(t,i,j)$]
    {$res\igobs letra(t,i,j)$}
    [$O(1)$]
    [Obtiene el momento en que una ficha del tablero fue puesta dada una coordenada $(i,j)$.]
    [\falta]
  \end{operaciones}
\end{interfaz}
\subsubsection{Implementación}
\subsubsection*{Representación}
\begin{Estructura}{tab}[tab\_estr]
  \textbf{donde} \texttt{tab\_estr} es \texttt{array\_dimensionable(array\_dimensionable(puntero(tupla(letra,nat))))}
\end{Estructura}
\subsubsection*{Invariante de Representación}
\Rep[tab\_estr]{$
  (\forall i:\texttt{nat})(i<tam(tab\_estr)\impluego tam(tab\_estr[i])=tam(tab\_estr))
$}

\subsubsection*{Función de Abstracción}
\Abs[tab\_estr]{tab}[e]{T}{
  $
  tam(e)\igobs tama\tilde{n}o(T)\yluego\\
  (\forall i,j:\texttt{nat})((enTablero?(T,i,j)\yluego hayLetra?(T,i,j))\impluego
  $\begin{adjustwidth}{12em}{0em}$
  (e[i][j]\neq NULL\yluego letra(T,i,j)\igobs\pi_{1}(*e[i][j])))
  $\end{adjustwidth}$
  $
}

\subsubsection*{Algoritmos}
\begin{algorithm}[H]
  \defun{iNuevoTablero}{\In{n}{nat}}{tab\_estr}
  \begin{algorithmic}[1]
    \State\asignar{filas}{\textsc{CrearArreglo}(n)}
    \For{$columnas\in filas$}\Comment{$O(N)$}
        \State\asignar{columnas}{\textsc{CrearArreglo}(n)}
        \For{$celda\in columnas$}\Comment{$O(N)$}
          \State\asignar{celda}{\textsc{NULL}}
        \EndFor
    \EndFor
    \State\Return{res}
  \end{algorithmic}
\end{algorithm}

\begin{algorithm}[H]
  \defun{iTamaño}{\In{t}{tab\_estr}}{nat}
  \begin{algorithmic}[1]
    \State\Return{$tam(t)$}
  \end{algorithmic}
\end{algorithm}

\begin{algorithm}[H]
  \defunv{iPonerLetra}{\Inout{t}{tab\_estr}, \In{i}{nat}, \In{j}{nat}, \In{l}{letra}, \In{tm}{nat}}
  \begin{algorithmic}[1]
    \State\asignar{t[i][j]}{\&\langle l,tm\rangle}
  \end{algorithmic}
\end{algorithm}

\begin{algorithm}[H]
  \defunv{iSacarLetra}{\Inout{t}{tab\_estr}, \In{i}{nat}, \In{j}{nat}}
  \begin{algorithmic}[1]
    \State\asignar{t[i][j]}{\textsc{NULL}}
  \end{algorithmic}
\end{algorithm}

\begin{algorithm}[H]
  \defun{iEnTablero}{\In{t}{tab\_estr}, \In{i}{nat}, \In{j}{nat}}{bool}
  \begin{algorithmic}[1]
    \State\Return{$i<tam(t)\land j<tam(t)$}
  \end{algorithmic}
\end{algorithm}

\begin{algorithm}[H]
  \defun{iHayLetra}{\In{t}{tab\_estr}, \In{i}{nat}, \In{j}{nat}}{bool}
  \begin{algorithmic}[1]
    \State\Return{$t[i][j]\neq\textsc{NULL}$}
  \end{algorithmic}
\end{algorithm}

\begin{algorithm}[H]
  \defun{iLetra}{\In{t}{tab\_estr}, \In{i}{nat}, \In{j}{nat}}{letra}
  \begin{algorithmic}[1]
    \State\Return{$\pi_{1}(*t[i][j])$}
  \end{algorithmic}
\end{algorithm}

\begin{algorithm}[H]
  \defun{iTiempo}{\In{t}{tab\_estr}, \In{i}{nat}, \In{j}{nat}}{nat}
  \begin{algorithmic}[1]
    \State\Return{$\pi_{2}(*t[i][j])$}
  \end{algorithmic}
\end{algorithm}

\subsection{Módulo Ocurrencia}
Es renombre de \texttt{conj(tupla(nat,nat,letra))} con las siguientes operaciones auxiliares.\\\\
\begin{operaciones}
  \InterfazFuncion{esHorizontal?}
  {
    \In{o}{ocurrencia}
  }{bool}
  [true]
  {$res\igobs true\iff(\forall f,f':\texttt{tupla(nat,nat,letra)})(f,f'\in o\land f\neq f'\implies \pi_{2}(f)=\pi_{2}(f'))$}
  [$O(\#o^{2})$]
  [Determina si una ocurrencia está alineada horizontalmente.]
  [\falta]\\

  \noindent\InterfazFuncion{esVertical?}
  {
    \In{o}{ocurrencia}
  }{bool}
  [true]
  {$res\igobs true\iff(\forall f,f':\texttt{tupla(nat,nat,letra)})(f,f'\in o\land f\neq f'\implies \pi_{1}(f)=\pi_{1}(f'))$}
  [$O(\#o^{2})$]
  [Determina si una ocurrencia está alineada verticalmente.]
  [\falta]\\

  \noindent\InterfazFuncion{haySuperpuestas?}
  {
    \In{o}{ocurrencia}
  }{bool}
  [true]
  {$res\igobs false\iff
      (\forall p,q:\texttt{nat})(\forall l,l':\texttt{letra})(
      \{\langle p,q,l\rangle,\langle p,q,l'\rangle\}\subseteq o\implies l=l')$}
  [$O(\#o^{2})$]
  [Determina si existen fichas distintas en una misma coordenada.]
  [\falta]
\end{operaciones}

\subsubsection*{Algoritmos}
\begin{algorithm}[H]
  \defun{iEsHorizontal?}{\In{o}{ocurrencia}}{bool}
  \begin{algorithmic}[1]
    \State\asignar{colns}{\textsc{Vacío}()}\Comment{Conjunto Lineal}
    \For{$ $\asignar{oIt}{\textsc{CrearIt}(o)}$;\ \textsc{HaySiguiente}(oIt);\ \textsc{Avanzar}(oIt)$}\Comment{$O(\#o)$}
      \State\asignar{ficha}{\textsc{Siguiente}(oIt)}
      \State$\textsc{Agregar}(colns,\pi_{2}(ficha)))$\Comment{$O(\#o)$}
    \EndFor
    \State\Return{\textsc{Cardinal}(colns)=1}
  \end{algorithmic}
\end{algorithm}

\begin{algorithm}[H]
  \defun{iEsVertical?}{\In{o}{ocurrencia}}{bool}
  \begin{algorithmic}[1]
    \State\asignar{filas}{\textsc{Vacío}()}\Comment{Conjunto Lineal}
    \For{$ $\asignar{oIt}{\textsc{CrearIt}(o)}$;\ \textsc{HaySiguiente}(oIt);\ \textsc{Avanzar}(oIt)$}\Comment{$O(\#o)$}
      \State\asignar{ficha}{\textsc{Siguiente}(oIt)}
      \State$\textsc{Agregar}(filas,\pi_{1}(ficha)))$\Comment{$O(\#o)$}
    \EndFor
    \State\Return{\textsc{Cardinal}(filas)=1}
  \end{algorithmic}
\end{algorithm}

\begin{algorithm}[H]
  \defun{iHaySuperpuestas?}{\In{o}{ocurrencia}}{bool}
  \begin{algorithmic}[1]
    \For{$ $\asignar{oIt}{\textsc{CrearIt}(o)}$;\ \textsc{HaySiguiente}(oIt);\ \textsc{Avanzar}(oIt)$}\Comment{$O(\#o)$}
        \State\asignar{ficha}{\textsc{Siguiente}(oIt)}
        \For{$ $\asignar{oItSig}{\textsc{Avanzar}(copy(oIt))}$;\ \textsc{HaySiguiente}(oItSig);\ \textsc{Avanzar}(oItSig)$}\Comment{$O(\#o)$}
          \State\asignar{ficha'}{\textsc{Siguiente}(oItSig)}
          \If{$\pi_{1}(ficha)=\pi_{1}(fichaSig)\land\pi_{2}(ficha)=\pi_{2}(fichaSig)$}
          \State\Return true
          \EndIf
        \EndFor
    \EndFor
    \State\Return false
  \end{algorithmic}
\end{algorithm}

\subsection{Módulo Notificación}
Asumimos que existe el tipo notif

%%% Local Variables:
%%% mode: latex
%%% TeX-master: "main"
%%% End:


\end{document}
